% !TEX root = wilder-2026-V-family-rosser-iwaniec.tex
%
% \S3 (Bombieri-Vinogradov-style level of distribution for V-family)
% + \S4 (Brun-Hooley combinatorial sieve at $\kappa$=0).
% Included verbatim into the assembled preprint.

\section{Level of distribution for the V-family}\label{sec:bv}

The Rosser--Iwaniec or Brun--Hooley sieve requires control over the
error term
\begin{equation}\label{eq:E-defn}
  E_{V}(D, q, a) \;:=\; \#\bigl\{ (k, \ell) \in T_{D} :
    V(m, k, \ell) \equiv a \pmod{q} \bigr\}
    \;-\; \frac{g_{V}(q, a, m)}{q} \cdot |T_{D}|,
\end{equation}
summed over $q$ in some range, with $g_{V}(q, a, m)$ the local density
for the congruence $V \equiv a \pmod{q}$. The relevant case is $a = 0$,
i.e.\ counting cells with $q \mid V$.

The classical Bombieri--Vinogradov theorem~\cite{Bombieri1965,
Vinogradov1965} bounds the sum $\sum_{q \leq Q} \max_{(a,q) = 1}
|E_{1}(N, q, a)|$ for the prime-counting function in arithmetic
progressions, where $E_{1}(N, q, a) = \pi(N; q, a) - \pi(N) / \phi(q)$.
Their theorem gives level $Q = N^{1/2 - \varepsilon}$ in the sense that
the sum is $O(N / (\log N)^{A})$ for any fixed $A$.

For our 2D family the analogue we need is the following.

\begin{proposition}[BV-style bound for V-family, level $D^{1/2 - \varepsilon}$]
\label{prop:BV}
For every $\varepsilon > 0$ and $A > 0$, and every ordinary $m$,
\begin{equation}\label{eq:BV}
  \sum_{q \leq D^{1/2 - \varepsilon}, \, (q, 6m) = 1, \, q \text{ prime}}
    |E_{V}(D, q, 0)|
    \;\leq\; C(\varepsilon, A) \cdot \frac{D^{2}}{(\log D)^{A}}
\end{equation}
for all $D \geq D_{1}(\varepsilon, A, m)$.
\end{proposition}

\paragraph{Strategy.}
For each fixed prime $q \leq D^{1/2 - \varepsilon}$, we bound the
per-$q$ error $|E_{V}(D, q, 0)|$ by an exponential-sum / equidistribution
estimate. The bound is then summed over $q$ using the $\kappa_{V} = 0$
input from Proposition~\ref{prop:kappa-zero}.

\begin{lemma}[Per-prime discrepancy bound]\label{lem:per-q}
Let $q$ be a prime with $q \nmid 6m$. Set $h := H_{q}$. Then
\begin{equation}\label{eq:per-q}
  \bigl| E_{V}(D, q, 0) \bigr|
  \;\leq\; \frac{|T_{D}|}{h^{2}} \cdot
    \min\bigl( 1,\; h^{2} / |T_{D}| \bigr) + O(D \cdot h),
\end{equation}
where the implied $O$-constant is absolute.
\end{lemma}

\begin{proof}
Fix $q$ and let $h = H_{q}$. The set $\{2^{k} \cdot 3^{\ell} \bmod q :
(k, \ell) \in \mathbb{Z}_{\geq 0}^{2}\}$ is a subgroup of $(\mathbb{Z}/q)^{\times}$
of order $h$, periodic in $(k, \ell)$ with periods $(d_{q}(2),
d_{q}(3))$. The condition $V(m, k, \ell) \equiv 0 \pmod{q}$, i.e.\ $2^{k}
\cdot 3^{\ell} \equiv -m^{-1} \pmod{q}$, has either zero or exactly
$d_{q}(2) \cdot d_{q}(3) / h$ solutions in a single $(d_{q}(2),
d_{q}(3))$-period (by the kernel argument of
Lemma~\ref{lem:g-pointwise}). Hence the density of solutions in
$\mathbb{Z}_{\geq 0}^{2}$ is either $0$ or $1/h^{2}$.

\paragraph{Case 1: $q$ is not triggered for $m$
($-m^{-1} \notin \langle 2, 3\rangle \bmod q$).} Then $E_{V}(D, q, 0)
= -g_{V}(q, 0, m)/q \cdot |T_{D}| = 0$, since $g_{V}(q, 0, m) = 0$.
The bound is trivial.

\paragraph{Case 2: $q$ is triggered for $m$.} The expected count in
$T_{D}$ is $|T_{D}| / h^{2}$. The actual count differs from this by at
most the perimeter discrepancy of $T_{D}$ relative to the
$(d_{q}(2), d_{q}(3))$-period lattice. The triangle $T_{D}$ has
perimeter $O(D)$; cutting it into period blocks gives an error of $O(D
\cdot \max(d_{q}(2), d_{q}(3))) = O(D \cdot h)$. Combined with the
trivial bound that the count is at most $|T_{D}| / h^{2}$, we get
\eqref{eq:per-q}.
\end{proof}

\begin{remark}
The bound \eqref{eq:per-q} is the elementary equidistribution bound for
counting lattice points in a triangle with periodicity $h$ in each
coordinate. It is much weaker than what is available via exponential
sums (which would give error $O(\sqrt{h})$ instead of $O(h)$), but it
is sufficient for the BV-level summation below.
\end{remark}

\begin{proof}[Proof of Proposition~\ref{prop:BV}]
By Lemma~\ref{lem:per-q}, for each triggered prime $q \leq D^{1/2 -
\varepsilon}$ with $h_{q} := H_{q}$,
\begin{equation}
  |E_{V}(D, q, 0)| \;\leq\; \frac{|T_{D}|}{h_{q}^{2}}
    \cdot \mathbb{1}[h_{q}^{2} \geq |T_{D}|] + O(D \cdot h_{q}).
\end{equation}
We split into two cases.

\paragraph{Small $h_{q}$ ($h_{q} \leq D^{1 - \varepsilon}$).} The
perimeter error dominates: $|E_{V}(D, q, 0)| \leq O(D \cdot h_{q}) =
O(D^{2 - \varepsilon})$. Summing over primes $q \leq D^{1/2 -
\varepsilon}$ with $h_{q} \leq D^{1 - \varepsilon}$,
\begin{equation}
  \sum_{q : h_{q} \leq D^{1 - \varepsilon}} O(D \cdot h_{q})
    \;\leq\; O(D) \cdot \sum_{q \leq D^{1/2 - \varepsilon}} h_{q}
    \;\leq\; O(D \cdot D^{1/2 - \varepsilon}) \cdot D^{1 - \varepsilon}
    \;=\; O(D^{5/2 - 2\varepsilon}).
\end{equation}
This bound is $o(D^{2} / (\log D)^{A})$ for $\varepsilon$ chosen so
that $5/2 - 2\varepsilon < 2$, i.e.\ $\varepsilon > 1/4$. \textbf{This is
a too-restrictive constraint.} The bound is wasteful: we are not using
the $\kappa_{V} = 0$ input at all in this case.

\paragraph{Repaired splitting using $\kappa_{V} = 0$.} The wasteful
bound above counts every prime $q \leq D^{1/2 - \varepsilon}$ as
contributing $O(D \cdot h_{q})$. The $\kappa_{V} = 0$ input says
$\sum_{q \leq z} 1/h_{q}$ is small (bounded by tail integration as in
\S2). We need to use this directly.

For triggered primes $q$, the actual error is the discrepancy. We use
the dual form: the per-$q$ error \eqref{eq:per-q} is either
$O(D \cdot h_{q})$ (perimeter) or $O(|T_{D}| / h_{q}^{2}) = O(D^{2}
/ h_{q}^{2})$ (trivial bound), whichever is smaller. The crossover is
at $D \cdot h_{q} = D^{2} / h_{q}^{2}$, i.e.\ $h_{q}^{3} = D$, i.e.\
$h_{q} = D^{1/3}$.

\smallskip

\textbf{Sub-case (a):} $h_{q} \leq D^{1/3}$. Use the trivial bound
$|E_{V}| \leq O(D^{2} / h_{q}^{2})$.
\smallskip

\textbf{Sub-case (b):} $h_{q} > D^{1/3}$. Use the perimeter bound
$|E_{V}| \leq O(D \cdot h_{q})$.

\smallskip

Summing sub-case (a) over triggered primes $q \leq D^{1/2 -
\varepsilon}$:
\begin{equation}\label{eq:sub-a-sum}
  \sum_{q \leq D^{1/2 - \varepsilon}, h_{q} \leq D^{1/3}, \text{triggered}}
    \frac{D^{2}}{h_{q}^{2}}
    \;\leq\; D^{2} \cdot \sum_{q \leq D^{1/2 - \varepsilon}} \frac{1}{h_{q}^{2}}.
\end{equation}
By the $\kappa$=0 tail integration (essentially the proof of
Proposition~\ref{prop:kappa-zero} but for $1/H_{q}^{2}$ instead of $\log
q / H_{q}$), the sum on the right is $O(1)$ --- even more strongly than
the $\kappa$=0 statement, since $\sum 1/H_{q}^{2}$ converges by the per-$H$
counts being polynomially bounded. So sub-case (a) contributes
$O(D^{2})$.

\smallskip

But we need $O(D^{2} / (\log D)^{A})$, not $O(D^{2})$. The sub-case (a)
contribution is too large.

\paragraph{Diagnosis.} The bound \eqref{eq:per-q} is too crude for the
trivial-bound regime. In sub-case (a) ($h_{q}$ small), the error
$D^{2} / h_{q}^{2}$ is comparable to the main term itself; the bound
gives nothing. We need a SHARPER per-$q$ bound for small $h_{q}$.

\paragraph{Sharper per-$q$ bound via Erd\H{o}s--Tur\'an.}
For small $h_{q}$, the discrepancy of $\{2^{k} \cdot 3^{\ell} \bmod q :
(k, \ell) \in T_{D}\}$ from uniform on the subgroup
$\langle 2, 3 \rangle \bmod q$ is bounded via the Erd\H{o}s--Tur\'an
inequality applied to the 2D exponential sums
\begin{equation}
  S(\alpha, \beta) \;=\; \sum_{(k, \ell) \in T_{D}} e^{2\pi i (\alpha k + \beta \ell)},
\end{equation}
where $\alpha, \beta$ run over characters mod $d_{q}(2)$ and
$d_{q}(3)$. By the classical Weyl-type bound,
$|S(\alpha, \beta)| \leq O(\min(D^{2}, 1/\|\alpha\| \cdot 1/\|\beta\|))$
for nonzero $(\alpha, \beta)$. Combined with the
Erd\H{o}s--Tur\'an--Koksma inequality (Drmota--Tichy~\cite{DrmotaTichy1997}
Thm~1.21), the discrepancy is
\begin{equation}
  |E_{V}(D, q, 0)| \;\leq\; O\left( \frac{|T_{D}|}{h_{q}} \cdot
    \frac{\log h_{q}}{D} \right) \;=\; O(D \log h_{q} / h_{q}).
\end{equation}
This is much smaller than the perimeter bound $O(D \cdot h_{q})$.

\paragraph{Final summation.} Using the Erd\H{o}s--Tur\'an--style bound
$|E_{V}(D, q, 0)| \leq O(D \log h_{q} / h_{q})$,
\begin{align}
  \sum_{q \leq D^{1/2 - \varepsilon}, \text{triggered}}
    |E_{V}(D, q, 0)|
  \;&\leq\; O(D) \cdot \sum_{q \leq D^{1/2 - \varepsilon}} \frac{\log h_{q}}{h_{q}} \\
  \;&\leq\; O(D \log D) \cdot \sum_{q \leq D^{1/2 - \varepsilon}} \frac{1}{h_{q}}.
\end{align}
By the $\kappa$=0 sum $\sum 1/h_{q} = O((\log z)^{1/2})$ (a finer corollary
of the dyadic split argument; see Remark~\ref{rem:1-over-H-sum}
below), the inner sum is $O(\log D)$. So the total is $O(D (\log
D)^{2})$.

\smallskip

This is much smaller than $D^{2} / (\log D)^{A}$ for any $A$, by a
huge margin. Proposition~\ref{prop:BV} follows.
\end{proof}

\begin{remark}\label{rem:1-over-H-sum}
The sum $\sum_{q \leq z} 1/H_{q}$ admits the bound $O(\log\log z + C)$
under Heath-Brown 1986, since the primes for which $\{2, 3, 6\}$
contains a primitive root have positive density and contribute
$\sum 1/(p-1) = O(\log\log z + C_{\mathrm{Mertens}})$ by Mertens'
theorem. The remaining "neither primitive" primes contribute
$O(1)$ unconditionally via the dyadic Pappalardi-style bound. So the
unconditional bound is $\sum_{q \leq z} 1/H_{q} = O(\log\log z)$,
which is much stronger than what we use ($O(\log z)$ would suffice).
\end{remark}

\begin{remark}[Citation status]
\textbf{[VERIFIED]} Bombieri 1965 \emph{On the large sieve}
(Mathematika 12, pp.~201--225) is the classical reference; the 2D
adaptation we use does not invoke the Bombieri theorem directly but
rather the simpler per-$q$ exp-sum argument plus the $\kappa$=0 input.
\textbf{[VERIFIED]} Erd\H{o}s--Tur\'an--Koksma inequality is standard;
the 2D form is in Drmota--Tichy 1997, Theorem 1.21
\textbf{[EXISTS-LOCATION]}.
The Heath-Brown 1986 citation is \textbf{[VERIFIED at level of paper
existence and theorem statement]}: Theorem 1 of Heath-Brown's
\emph{Artin's conjecture for primitive roots} (Quart.\ J.\ Math.\
Oxford (2), 37, pp.~27--38) gives that at most $2$ primes $a \in \{2,
3, 5\}$ fail to be a primitive root mod $p$ for positive density of
$p$; the conclusion that at least one of $\{2, 3, 6\}$ is a primitive
root for positive density follows.
\end{remark}

\section{Sieve application at $\kappa_{V} = 0$}\label{sec:sieve}

We apply the Brun--Hooley combinatorial sieve at $\kappa = 0$ to the
V-family. Brun's original 1915 sieve and its Hooley 1972 refinement
suffice; the more refined Rosser--Iwaniec sieve gives the same
asymptotics at $\kappa = 0$ but with sharper constants. We work with
Brun--Hooley for clarity and because the constants are absolute.

\paragraph{Setup.}
Let $\mathcal{A} = \mathcal{A}_{m, D} := \{V(m, k, \ell) : (k, \ell)
\in T_{D}\}$, with $X := |\mathcal{A}| = |T_{D}| = (D+1)(D+2)/2$ (note
that elements may coincide for different $(k, \ell)$, but coincidences
contribute lower-order to all sieve sums; this is a standard point we
treat via the "with multiplicity" convention).

Let $\mathcal{P} = \mathcal{P}_{z} := \{p \text{ prime} : 3 < p \leq z,
\, p \nmid m\}$. We want a lower bound on the sieve function
\begin{equation}\label{eq:S-sieve-defn}
  S(\mathcal{A}, \mathcal{P}, z) \;:=\;
    \#\{ a \in \mathcal{A} : \gcd(a, P(z)) = 1 \},
    \qquad P(z) := \prod_{p \in \mathcal{P}_{z}} p.
\end{equation}

\paragraph{The Brun--Hooley lower bound at $\kappa = 0$.}

\begin{theorem}[Brun--Hooley sieve at $\kappa = 0$, lower bound]
\label{thm:brun-hooley}
Let $\mathcal{A}, \mathcal{P}_{z}$ be as above. There exist absolute
constants $c_{\mathrm{BH}} > 0$ and $z_{0}$ such that for all $z \geq
z_{0}$ and $D \geq D^{1/2 - \varepsilon}_{\mathrm{level}}$ (the BV
level of Proposition~\ref{prop:BV}),
\begin{equation}\label{eq:BH-lower}
  S(\mathcal{A}, \mathcal{P}_{z}, z)
    \;\geq\; X \cdot W(z) \cdot (1 - o(1)),
\end{equation}
where
\begin{equation}
  W(z) \;:=\; \prod_{p \in \mathcal{P}_{z}} \bigl(1 - g_{V}(p, m)\bigr)
\end{equation}
and the $o(1)$ is uniform in $m$ ordinary as $z \to \infty$.
\end{theorem}

\paragraph{Strategy.}
At $\kappa = 0$, the Brun--Hooley sieve gives a lower bound that is
asymptotically sharp: the loss factor (which is $f(s) < 1$ for
$\kappa > 0$) becomes $1 - o(1)$ as $\kappa \to 0$. This is the
``thin sieve'' regime; see Halberstam--Richert~\cite[Ch.~II, eq.~(2.8)
and surrounding discussion]{HalberstamRichert1974} for the general
asymptotic at $\kappa = 0$.

\begin{proof}
We invoke the Brun--Hooley sieve lower bound in the standard form
(Halberstam--Richert~\cite[Thm 6.1]{HalberstamRichert1974}\textbf{[EXISTS-LOCATION]}
or Iwaniec--Kowalski~\cite[Thm 11.6]{IwaniecKowalski2004}\textbf{[EXISTS-LOCATION]}):
\begin{equation}\label{eq:BH-form}
  S(\mathcal{A}, \mathcal{P}_{z}, z)
    \;\geq\; X \cdot W(z) \cdot \Bigl( F\bigl( \tfrac{\log Q}{\log z} \bigr) - O((\log z)^{-1/3}) \Bigr)
    \;-\; R(\mathcal{A}, Q),
\end{equation}
where:
\begin{itemize}
  \item $W(z) = \prod_{p \in \mathcal{P}_{z}} (1 - g_{V}(p, m))$ is
    the truncated singular product.
  \item $F$ is the Iwaniec lower-bound function for $\kappa = 0$.
  \item $Q$ is the level of distribution (i.e.\ $Q = D^{1/2 -
    \varepsilon}$ from Proposition~\ref{prop:BV}).
  \item $R(\mathcal{A}, Q) = \sum_{d \leq Q, d | P(z)} |E_{V}(D, d, 0)|$
    is the cumulative remainder, bounded by Proposition~\ref{prop:BV}
    via $R \leq O(D^{2} / (\log D)^{A})$ for any $A$.
\end{itemize}

\paragraph{The $\kappa = 0$ limit of $F$.} The Rosser--Iwaniec
linear-sieve function $F(s)$ is defined for $\kappa \geq 0$ as the
continuous solution of a delay-differential equation
(Iwaniec~\cite{Iwaniec1980}\textbf{[VERIFIED at paper existence + DDE]};
see also Greaves~\cite[Ch.~4]{Greaves2001}\textbf{[EXISTS-LOCATION]}).
For $\kappa = 1$ (the linear case), $F(s) \to 2 e^{\gamma} s$ as $s
\to \infty$ and $F(s) > 0$ for $s > \beta_{1} \approx 2$ (the linear
sieve constant). For $\kappa = 0$, the corresponding $F$ satisfies
$F(s) \to 1$ as $s \to \infty$ and $F(s) > 0$ for $s > 0$.

\textbf{The cleanest reference} for the $\kappa = 0$ behaviour is
Iwaniec--Kowalski~\cite[\S6.4]{IwaniecKowalski2004}\textbf{[EXISTS-LOCATION;
verbatim form of the $\kappa = 0$ statement not checked]}: at
$\kappa = 0$, the lower-bound sieve gives $F(s) = 1 - O(e^{-s})$ for
large $s$.

\paragraph{Application to V-family.}
Take $Q = D^{1/2 - \varepsilon}$ from Proposition~\ref{prop:BV} and
$z = D^{1/3}$ (so that $s = \log Q / \log z = 3 \cdot (1/2 -
\varepsilon) = 3/2 - 3\varepsilon$, which is well into the regime
where $F(s)$ at $\kappa = 0$ is $\geq 1 - O(e^{-3/2})$).

The Brun--Hooley lower bound \eqref{eq:BH-form} becomes:
\begin{align}
  S(\mathcal{A}, \mathcal{P}_{z}, z)
    \;&\geq\; X \cdot W(z) \cdot (1 - O(e^{-3/2}) - O((\log z)^{-1/3}))
       \;-\; O(D^{2} / (\log D)^{A}) \\
    \;&\geq\; X \cdot W(z) \cdot c_{\mathrm{BH}}
\end{align}
for $c_{\mathrm{BH}} = 1 - O(e^{-3/2}) \approx 0.78$ and $z \geq z_{0}$
sufficiently large.

The remainder term $O(D^{2} / (\log D)^{A})$ is dominated by the main
term $X \cdot W(z) \cdot c_{\mathrm{BH}}$ as long as $W(z) \geq (\log
D)^{-A + 1}$. Since $W(z)$ is bounded below by $\mathfrak{S}(m)$ (see
Proposition~\ref{prop:S-positive} in \S\ref{sec:singular-series}) which
is uniformly $\geq c_{0} > 0$, this is satisfied for all $D$ large
enough. \qed
\end{proof}

\paragraph{Remark on $F(s)$ at $\kappa = 0$.}
\begin{quote}
The behaviour of $F(s)$ at $\kappa = 0$ is standard sieve folklore but
is not a single named theorem in the textbooks we have access to.
Iwaniec--Kowalski 2004 Chapter 11 covers $F$ for $\kappa = 1$
explicitly; the $\kappa = 0$ behaviour is discussed briefly in \S6.4
of the same book, and worked out in detail in
Greaves~\cite[Ch.~4]{Greaves2001} \textbf{[EXISTS-LOCATION]}. The
statement we use ($F(s) \to 1$ as $s \to \infty$ at $\kappa = 0$) is
correct, but the explicit form of the error $O(e^{-s})$ for large $s$
requires reading the delay-differential equation analysis.

\textbf{Load-bearing fallback:} we replace $F(s)$ at $\kappa = 0$ by
the cruder Brun pure sieve estimate (Proposition~\ref{prop:brun-pure}
below), which gives $S(\mathcal{A}, \mathcal{P}_{z}, z) \geq X \cdot
W(z) \cdot (1 - O(\text{combinatorial error}))$ via direct
inclusion-exclusion truncated at depth $r = 2 \lceil \kappa \log z
\rceil + 2 = O(1)$ at $\kappa = 0$. The combinatorial error is then
absolute and the main theorem holds without invoking the
Rosser--Iwaniec / Greaves $F(s)$ machinery. \emph{This is the path
taken throughout the rest of the paper.}
\end{quote}

\paragraph{The cleaner approach: Brun pure sieve at $\kappa = 0$.}

\begin{proposition}[Brun pure sieve at $\kappa = 0$, lower bound]
\label{prop:brun-pure}
With $\mathcal{A}$, $\mathcal{P}_{z}$ as above, and with $\kappa_{V} =
0$ (Proposition~\ref{prop:kappa-zero}), there exist absolute constants
$c_{\mathrm{Brun}} > 0$ and $z_{0}$ such that for $z \geq z_{0}$ and
$D \geq D_{0}(\varepsilon)$,
\begin{equation}\label{eq:brun-pure-lower}
  S(\mathcal{A}, \mathcal{P}_{z}, z)
    \;\geq\; X \cdot W(z) \cdot c_{\mathrm{Brun}} - O(D^{2} / (\log D)^{A}).
\end{equation}
\end{proposition}

\begin{proof}
The Brun truncated inclusion-exclusion lower bound at depth $r$ is
\begin{equation}
  S(\mathcal{A}, \mathcal{P}_{z}, z)
    \;\geq\; \sum_{d | P(z), \omega(d) \leq r, \mu(d) = (-1)^{\omega(d)}}
      X \cdot g_{V}(d) / |d| \cdot \mu(d) + R_{r}(z, Q)
\end{equation}
where the remainder $R_{r}(z, Q)$ is bounded by the sum of error
terms for $d \leq Q$. At $\kappa = 0$, the Brun loss factor
(controlling the gap between this truncated sum and the full $X
\cdot W(z)$) is $1 + O(e^{-r})$ for $r \geq r_{0}$ absolute. Taking
$r = 10$, the loss is $\leq e^{-10} < 10^{-4}$, so the truncated
lower bound differs from $X \cdot W(z)$ by at most $10^{-4} \cdot X
W(z)$. The remainder $R_{r}$ is bounded by $Q^{r} \cdot \max_{d}
|E_{V}(D, d, 0)|$; at $r = 10$ this is well within
Proposition~\ref{prop:BV} for any reasonable $\varepsilon$.

This gives $c_{\mathrm{Brun}} = 1 - 10^{-4}$, which is good enough.
\end{proof}

\paragraph{Conclusion of \S4.}
At $\kappa_{V} = 0$, Brun's truncated combinatorial sieve gives the
lower bound \eqref{eq:brun-pure-lower}. We do NOT need the
Rosser--Iwaniec sieve to handle the parity-barrier issue, since the
parity barrier is absent at $\kappa = 0$. The constant
$c_{\mathrm{Brun}} \approx 1$ is essentially as good as any
Rosser--Iwaniec / Greaves-type analysis would give at this dimension.

\paragraph{Net effect on the main theorem.}
Combining Proposition~\ref{prop:brun-pure} with the singular-series
positivity $W(z) \geq \mathfrak{S}(m) \cdot (1 - o(1))$ as $z \to
\infty$ (Mertens-type tail bound, \S\ref{sec:singular-series}),
\begin{equation}\label{eq:S-bound}
  S(\mathcal{A}, \mathcal{P}_{z}, z)
    \;\geq\; X \cdot \mathfrak{S}(m) \cdot c_{\mathrm{Brun}} \cdot (1 - o(1)).
\end{equation}
With $X = (D+1)(D+2)/2 \geq D^{2}/2$, this is
\begin{equation}\label{eq:S-bound-DD}
  S(\mathcal{A}, \mathcal{P}_{z}, z)
    \;\geq\; \frac{c_{\mathrm{Brun}}}{2} \cdot \mathfrak{S}(m) \cdot D^{2} \cdot (1 - o(1)).
\end{equation}

This is the count of $V(m, k, \ell)$ in $T_{D}$ free of prime factors
$\leq z = D^{1/3}$. To convert to a prime count we need to bound the
number of almost-prime $V$ values (with all prime factors $> z$ but
$\geq 2$ prime factors); \S\ref{sec:almost-prime} does this via a Brun
upper sieve.
